\documentclass[12pt,a4paper]{report}
\usepackage[utf8]{inputenc}
\usepackage{amsmath}
\usepackage{amsfonts}
\usepackage{amssymb}
\usepackage{amsthm}
\usepackage{hyperref}

\usepackage{multicol}
\usepackage{fancyhdr}
\usepackage[inline]{enumitem}
\usepackage{tikz}
\usepackage{tikz-cd}
\usetikzlibrary{calc}
\usetikzlibrary{shapes.geometric}
\usetikzlibrary{positioning}
\usepackage[margin=0.5in]{geometry}
\usepackage{xcolor}

\hypersetup{
    colorlinks=true,
    linkcolor=blue,
    filecolor=magenta,      
    urlcolor=cyan,
    pdftitle={Tensors},
    pdfpagemode=FullScreen,
    }

%\urlstyle{same}

\newcommand{\CLASSNAME}{Math 5050 -- Special Topics: Differential Geometry}
\newcommand{\STUDENTNAME}{Paul Carmody}
\newcommand{\ASSIGNMENT}{Section 19: The Exterior Derivative }
\newcommand{\DUEDATE}{October 8, 2025}
\newcommand{\PROFESSOR}{Professor Berchenko-Kogan}
\newcommand{\SEMESTER}{Fall 2025}
\newcommand{\SCHEDULE}{TBD}
\newcommand{\ROOM}{Remote}

\newcommand{\MMN}{M_{m\times n}}
\newcommand{\FF}{\mathcal{F}}

\pagestyle{fancy}
\fancyhf{}
\chead{ \fancyplain{}{\CLASSNAME} }
%\chead{ \fancyplain{}{\STUDENTNAME} }
\rhead{\thepage}
\input{../MyMacros}

\newcommand{\RED}[1]{\textcolor{red}{#1}}
\newcommand{\BLUE}[1]{\textcolor{blue}{#1}}

\newcommand{\SUPP}{\operatorname{supp}}
\newcommand{\CLOSURE}{\operatorname{closure of}}

\begin{document}

\begin{center}
	\Large{\CLASSNAME -- \SEMESTER} \\
	\large{ w/\PROFESSOR}
\end{center}
\begin{center}
	\STUDENTNAME \\
	\ASSIGNMENT -- \DUEDATE\\
\end{center} 


\begin{enumerate}[label=19.\arabic*]

\item \textbf{Pullback of a differential form}

Let $U$ be the open set $(0,\infty)\times(0,\pi)\times(0,2\pi)$ in the $(\rho,\theta,\theta)$-space $\R^3$.  Define $F: U \to \R^3$ by
\begin{align*}
	F(\rho, \phi, \theta) =(\rho\sin\phi\cos\theta,\rho\sin\phi\sin\theta,\rho\cos\phi).
\end{align*}If $x,y,z$ are the standard coordinates on the target $\R^3$, show that 
\begin{align*}
	F^*(dx\wedge dy\wedge dz) = \rho^2\sin\phi\, d\rho\wedge d\phi\wedge d\theta.
\end{align*}

\BLUE{Start by recognizing that the $n$-form in an $n$-dimensional vector space operating on a function, $F: N \to M$ is the determinant of the Jacobian, $J_F$.
\begin{align*}
%	J_F &= \THREEXTHREE{\PART{F_x}{\rho}}{ \PART{F_x}{\phi}}{\PART{F_x}{\theta}}
%		{\PART{F_y}{\rho}}{ \PART{F_y}{\phi}}{\PART{F_y}{\theta}}
%		{\PART{F_z}{\rho}}{ \PART{F_z}{\phi}}{\PART{F_z}{\theta}}
	J_F &= \COLVECTOR{ dx \\ dy \\ dz } \\
	dx &= \CYCLE{\PART{F_x}{\rho} &  \PART{F_x}{\phi} & \PART{F_x}{\theta}} \\
		&= \CYCLE{\sin\phi\cos\theta & \rho\cos\phi\sin\theta & -\rho\sin\phi\sin\theta} \\
	dy &= \CYCLE{\PART{F_y}{\rho} &  \PART{F_y}{\phi} & \PART{F_y}{\theta}} \\
	&= \CYCLE{\sin\phi\cos\theta & \rho\cos\phi\sin\theta & -\rho\sin \phi\cos\theta }\\
	dz &= \CYCLE{\PART{F_z}{\rho} &  \PART{F_z}{\phi} & \PART{F_z}{\theta}} \\
	&= \CYCLE{\cos \phi & -\rho\sin\phi & 0} \\
	\BARS{J_F} &= \cos\phi\DTWOXTWO{\rho\cos\phi\sin\theta}{ -\rho\sin\phi\sin\theta}{ \rho\cos\phi\sin\theta}{ -\rho\sin \phi\cos\theta }+\rho\sin\theta\DTWOXTWO{\sin\phi\cos\theta }{ -\rho\sin\phi\sin\theta}{\sin\phi\cos\theta }{ -\rho\sin \phi\cos\theta}\\
	&= \rho^2\PAREN{ \cos\phi\DTWOXTWO{\cos\phi\sin\theta}{ -\sin\phi\sin\theta}{ \cos\phi\sin\theta}{ -\sin \phi\cos\theta }+\sin\phi\DTWOXTWO{\sin\phi\cos\theta }{ -\sin\phi\sin\theta}{\sin\phi\cos\theta }{ -\sin \phi\cos\theta}}\\
	&= \rho^2\PAREN{ \cos^2\phi\sin\phi \DTWOXTWO{\sin\theta}{ -\sin\theta}{ \sin\theta}{ -\cos\theta }+\sin^3\phi\DTWOXTWO{\cos\theta }{ -\sin\theta}{\cos\theta }{ -\cos\theta}}\\
	&= \rho^2\PAREN{ \cos^2\phi\sin\phi \sin\theta\DTWOXTWO{1}{ -\sin\theta}{ 1}{ -\cos\theta }+\sin^3\phi\cos\theta \DTWOXTWO{1 }{ -\sin\theta}{1 }{ -\cos\theta}}\\
	&= \rho^2\sin\phi\PAREN{ \cos^2\phi \sin\theta\DTWOXTWO{1}{ -\sin\theta}{ 1}{ -\cos\theta }+\sin^2\phi\cos\theta \DTWOXTWO{1 }{ -\sin\theta}{1 }{ -\cos\theta}}
\end{align*}The bracketed term simplies
\begin{align*}
	\cos^2\phi \sin\theta\DTWOXTWO{1}{ -\sin\theta}{ 1}{ -\cos\theta }+\sin^2\phi\cos\theta \DTWOXTWO{1 }{ -\sin\theta}{1 }{ -\cos\theta} &= (\cos^2\phi \sin\theta+\sin^2\phi\cos\theta )( -\cos\theta  +\sin\theta) \\
	&=\cos^2\phi \sin^2\theta +\sin^2\phi\cos^\theta \\
	&= 1
\end{align*}Hence, $|J_F| = \rho^2\sin\phi$ .
}

\item \textbf{Pullback of a differential form}

Let $F: \R^2 \to \R^2$ be given by 
\begin{align*}
	F(x,y) = (x^2+y^2,xy).
\end{align*}If $u,v$ are the standard coordinates on the target $\R^2$, compute $F^*(udu+vdv)$.

\BLUE{For clarity, 
\begin{align*}
	F(x,y) = (x^2+y^2)dx + (xy)dy.
\end{align*}Therefore, we can derive $du$ and $dv$ as
\begin{alignat*}{2}
	du &= \PART{F_x}{x}dx + \PART{F_x}{y}dy, & \quad 
	dv &= \PART{F_y}{x}dx + \PART{F_y}{y}dy\\
	&= \PART{x^2+y^2}{x}dx + \PART{x^2+y^2}{y}dy, & \quad
	&= \PART{xy}{x}dx + \PART{xy}{y}dy \\
	&= 2xdx+2ydy, & \quad
	&= ydx+xdy
\end{alignat*}Let $\omega(u,v) = udu+vdv$ then
\begin{align*}
	F^*\omega(x,y) &= \omega(F(x,y)) \\
	&= (x^2+y^2)(2xdx+2ydy)+(xy)(ydx+xdy) \\
	&= 2x^3dx+2x^2ydy+2xy^2dx+2y^3dy+xy^2dx+x^2ydy \\
	&= (2x^3+2xy^2+xy^2)dx+(2x^2y+2y^3+x^2y)dy \\
	&= (2x^3+3xy^2)dx+(3x^2y+2y^3)dy
\end{align*}
}

\item \textbf{Pullback of a differential form by a curve}

Let $\tau$ be the 1-form $\tau = (-ydx+xdy)/(x^2+y^2)$ on $\R^2-\{\textbf{0}\}$.  Define $\gamma: \R \to \R^2 - \{\textbf{0}\}$ by $\gamma(t)=(\cos t,\sin t)$.  Compute $\gamma^*\tau$.  (this problem is related to Example 17.16 in that if $i: S^1\hookrightarrow \R^2-\{\textbf{0}\}$ is the inclusion, then $\gamma= i \circ c$ and $\omega=i^*\tau$).

\BLUE{\begin{align*}
	\tau &= \frac{-y}{x^2+y^2}dx+\frac{x}{x^2+y^2} dy \\
		&= \frac{1}{x^2+y^2}\PAREN{-ydx+xdy} \\
	\gamma^*\tau(t) &=\gamma^*\PAREN{\frac{1}{x^2+y^2}}\PAREN{-\gamma^*(y)d(\gamma^*x)+\gamma^*(x)d(\gamma^*y)} \\
	&= \PAREN{\frac{1}{\cos^2t+\sin^2 t}}\PAREN{-\sin t(-\sin t)dt+\cos t(\cos t)dt} \\
	 &= 1
\end{align*}
}

\item \textbf{Pullback of a restriction}

Let $F: N\to M$ be a $C^\infty$ map of manifolds, $U$ an open subset of $M$, and $\RBAR{F}_{F^{-1}(U)} :F^{-1}(U) \to U$ the restriction of $F$ to $F^{-1}(U)$.  Prove that if $\omega\in\Omega^k(M)$, then
\begin{align*}
	\PAREN{\RBAR{F}_{F^{-1}(U)}}^*(\omega|_U) = (F^*\omega)|_{F^{-1}(U)}.
\end{align*}

\item \textbf{Coordinate functions are differential forms}

Let $f^1,\dots,f^n$ be $C^\infty$ functions on a neighborhood $U$ of a point $p$ in a manifold of dimension $n$.  Show that there is a neighborhood $W$ of $p$ on which $f^1,\dots,f^n$ form a coordinate system if and only if $(df^1\wedge\dots\wedge df^n)_p \ne 0$.

\item \textbf{Local operators}

An opeerator $L:\Omega^*(M) \to \Omega^*(M)$ is \textit{support-decreasing} if $\SUPP L(\omega) \subset \SUPP\omega$ for every $k$-form $\omega \in \Omega^*(M)$ for all $k \ge 0$.  Show that an operator on $\Omega^*(M)$ is local if and only if it is support-decreasing.

\item \textbf{Derivations on $C^\infty$ functions are local operators}

Let $M$ be a smooth manifold.  The definition of a \textit{local operator}  $D$ on $C^\infty(M)$ is similar to that of a local operator on $\Omega^*(M): D$ is \textit{local} if whenever a function $f \in C^\infty(M)$ vanishes identically on an open subset $U$, then $Df \equiv 0$ on $U$.  Prove that a derivation of $C^\infty(M)$ is a local operator $C^\infty(M)$.

\item \textbf{Nondegenerate 2-forms}

Let 2-covector $\alpha$ on a 2n-dimensional vector space $V$ is said to be \textit{nondegenerate} if $\alpha^n:= \alpha\wedge\cdots\alpha$ ($n$ times) is not the zero 2n-covector.  A 2-form $\omega$ on a 2n-dimensional manifold $M$ is said to be \textit{nondegenerate} if at every point $p\in  M$, the 2-covector $\omega_p$ is nondegenerate on the tangent space $T_pM$.
\begin{enumerate}[label=(\alph*)]

	\item Prove that on $\C^n$ with real coordinates $x^1,y^1,\dots,x^n,y^n$, the 2-form 
	\begin{align*}
		\omega = \sum_{j=1}^n dx^j\wedge dy^j
	\end{align*}is degenerate.
	\item Prove that if $\lambda$ is the Liouville form on the total space $T^*M$ of the cotangent bundle of an $n$-dimensional manifold $M$, then $d\lambda$ is a nondegenerate 2-form on $T^*M$.

\end{enumerate}

\item \textbf{Vertical planes}

Let $x,y,z$ be the standard coordiantes on $\R^3$.  A plane in $\R^3$ is 	\textit{vertical} if it is defined by $ax+by=0$ for some $(a,b) \ne (0,0) \in \R^2$.  Prove that restricted to a vertical plane, $dx \wedge dy =0$.

\item \textbf{Nowhere-vanishing form on $S^1$}

Prove that the nowhere-vanishing form $\omega$ on $S^1$ constructed in Example 19.8 is the form $-ydx+xdy$ of Example 17.15.  (\textit{Hint:} Consier $U_x$ and $U_y$ separeately.  On $U_x$, substitute $dx=-(y/x)dy$ into $-ydx+xdy$.)

\item \textbf{A $C^\infty$ nowhere-vanishing form on a smooth hypersurface}

\begin{enumerate}[label=(\alph*)]

	\item Let $f(x,y)$ be a $C^\infty$ function on $\R^2$ and assume that 0 is a regular value of $f$.  Bey the regular level set theorem, the zero set $M$ of $f(x,y)$ is a one-dimensional submanifold of $\R^2$. Construct a $C^\infty$ nowhere-vanishing 1-form on $M$.
	
	\item Let $f(x,y,z)$ be a $C^\infty$ function on $\R^3$ and assume that 0 is a regular value $f$.  By the regular level set theorem, the zero set $M$ of $f(x,y,z)$ is a two-dimensional submanifold of $\R^3$.  Let $f_x, f_y, f_z$ be the partial derivatives of $f$ with respect to $x,y,z$ respectively.  Show that the equalities
	\begin{align*}
		\frac{dx\wedge dy}{f_z}=\frac{dy\wedge dz}{f_x}=\frac{dz\wedge dx}{f_y}
	\end{align*}hold on $M$ whenever they make sense, nad therfore the three 2-forms piece together to give $C^\infty$ nowhere-vanishing 2-form on $M$.
	
	\item Generalize this problem to the regular level set of $f(x^1,\dots,x^{n+1})$ in $\R^{n+1}$.

\end{enumerate}

\item \textbf{Vector fields as derivations of $C^\infty$ fucntions}

In Subsection 14.1 we showed that $C^\infty$ vector field $X$ on a manifold $M$ gives rise to a derivation of $C^\infty(M)$. We will now show that every derivation of $C^\infty(M)$ arise from one and only one vector field, as promised earlier. To distingisg the vector field from the derivation, we will temporarily denote the derivation arising from $X$ by $\varphi(X)$.  Thus, for any $f\in C^\infty(M)$,
\begin{align*}
	(\varphi(X)f)(p)=x_pf,\, \text{for all }p \in M.
\end{align*}\begin{enumerate}[label=(\alph*)]
	\item Let $\mathcal{F}= C^\infty(M)$.  Prove that $\varphi: \XXX(M) \to \DER(C^\infty(M))$ is an $\mathcal{F}$-linear map.
	
	\item Show that $\varphi$ is injective.
	\item If $D$ is a derivation of $C^\infty(M)$ and $p\in M$, defined $D_p:C_p^\infty(M) \to C_p^\infty(M)$ by
	\begin{align*}
		D_p[f] = [D\hat{f}] \in C_p^\infty(M),
	\end{align*}where $[f]$ is the germ of $f$ at $p$ and $\hat{f}$ is a global extension of $f$, such as those given by Proposition 18.8.  Show that $D_p[f]$ is well defined. (\textit{Hint:} Apply probloem 19.7.)
	
	\item Show that $D_p$ is a derivation of $C_p^\infty(M)$.
	\item Prove that $\varphi:\XXX(M)\to \DER(C^\infty(M))$ is an isomorphism of $\mathcal{F}$-modules.

\end{enumerate}

\item \textbf{Twentieth-century formulation of Maxwell's equations}

In Maxwell's theory of electricity and magnetism, developed in the late nineteenth century,  the electric field $\textbf{E}=\ABRACKET{E_1,E_2,E_3}$ and the magentic field $\textbf{B}=\ABRACKET{B_1,B_2,B_3}$ in a vacuum $\R^3$ with no charge or current satisfy the following equations.
\begin{alignat*}{2}
	\nabla\times\textbf{E}&= -\frac{\partial B}{\partial t}, & \quad \nabla\times B&=\frac{\partial E}{\partial t},\\
	\operatorname{div} E&=0, &\quad \operatorname{div}v B&=0
\end{alignat*}By the correspondence in Subsection 4.6, the 1-form $E$ on $\R^3$ corresponding to the vector field $E$ is 
\begin{align*}
	E=E_1dx+E_2dy+E_3dz
\end{align*}and the 2-form $B$ on $\R^3$ corresponding to the vector field $B$ is
\begin{align*}
	B=B_1 dy\wedge dz + B_2 dz\wedge dx + B_3 dx \wedge dy.
\end{align*}

Let $\R^4$ be space-time with coordinates $(x,y,z,t)$.  Then both $E$ and $B$ can be viewed as differential forms on $\R^4$. Define $F$ to be the 2-form
\begin{align*}
	F=E \wedge dt +B
\end{align*}on space-time.  Decide which two Maxwell's equations are equivalent to the equation
\begin{align*}
	dF=0.
\end{align*}Prove your answer. (The other tow are equivalent to $D*F=0$ for a star-operator $*$ defined in differential geometry.  See [2,Section 19.1, p. 689].)
\end{enumerate}

\end{document}
